% ============================================================
% Chapter 1: Introduction
% File: chapters/chapter1_introduction.tex
% ============================================================


\clearpage
\section*{Chapter 1}
\section{Introduction}
\label{sec:introduction}


% =========================================================================
% 1.1  MOTIVATION
% =========================================================================
\subsection{Hedging Under Non-Uniqueness of Martingale Measures}
\label{subsec:hedging_motivation}

The modern theory of derivative pricing rests on the principle that, in a
frictionless market with continuous trading, the price of a contingent
claim is determined by an arbitrage argument rather than by the
preferences of individual agents.  In the setting of
\citet{BlackScholes1973} and \citet{Merton1973}, a single risky asset
follows geometric Brownian motion with constant volatility~$\sigma$, and a
risk-free bond grows at a constant rate~$r$.  Under these assumptions,
every European option payoff can be exactly replicated by dynamically
adjusting the stock position according to an explicit formula, the
celebrated delta hedge.  The cost of constructing the replicating
portfolio at any given time is unique: it equals the discounted expectation
of the payoff under the sole probability measure that makes discounted
prices into martingales.

The mathematical foundation of this result is provided by the
{fundamental theorems of asset pricing}.  The first theorem,
established in full generality by \citet{DelbaenSchachermayer1994},
states that a market satisfies {no free lunch with vanishing risk}
(NFLVR) if and only if there exists at least one probability measure
$\mathbb{Q}$, equivalent to the physical measure~$\mathbb{P}$, under
which discounted asset prices are local martingales.  Any such measure is
called an {equivalent local martingale measure}, and the set of all
such measures is denoted $\mathcal{M}(X)$.  The second theorem, due to
\citet{HarrisonPliska1981} and \citet{HarrisonKreps1979}, characterises
market completeness: every contingent claim can be replicated by a
self-financing trading strategy if and only if $\mathcal{M}(X)$ contains
exactly one element.  When $|\mathcal{M}(X)| = 1$, the unique
measure~$\mathbb{Q}$ assigns a unique price to every claim, and the
{martingale representation theorem} guarantees the existence of a
replicating strategy whose gains-from-trade process equals the conditional
expectation of the payoff.

The structural issue that motivates this thesis is the failure of
uniqueness.  In practice, the assumptions underpinning the complete-market
paradigm are systematically violated.  Volatility is not constant: it
evolves randomly over time, as documented by the implied volatility
surface observed in options markets and formalised in stochastic volatility
models such as that of \citet{Heston1993}.  Asset prices exhibit jumps
driven by macroeconomic announcements or liquidity shocks
\citep{Merton1976}.  Transaction costs prevent the continuous rebalancing
that perfect replication requires \citep{LelandReplication1985}.  In
multi-asset settings, portfolios may be exposed to correlated risk factors
that are not individually tradable.  Each of these features introduces
sources of randomness that cannot be offset by the available instruments,
rendering the market {incomplete}.

When the market is incomplete, $\mathcal{M}(X)$ contains more than one
element.  Different measures in $\mathcal{M}(X)$ assign different prices
to the same non-replicable claim, and no single trading strategy can
eliminate the hedging error.  Rather than a unique price, one obtains an
interval of arbitrage-consistent values.  Let $H \ge 0$ denote the
discounted payoff of a European claim at maturity~$T$.  The
{arbitrage-free price interval} is
\begin{equation}
\label{eq:price_interval}
\Bigl[\,
  \inf_{\mathbb{Q} \in \mathcal{M}(X)} \mathbb{E}^{\mathbb{Q}}[H],
  \;\;
  \sup_{\mathbb{Q} \in \mathcal{M}(X)} \mathbb{E}^{\mathbb{Q}}[H]
\,\Bigr].
\end{equation}
Any price within this interval is consistent with the absence of
arbitrage; any price outside it creates a free lunch
\citep{DelbaenSchachermayer1994}.  The {upper bound},
$\pi(H) = \sup_{\mathbb{Q} \in \mathcal{M}(X)}
\mathbb{E}^{\mathbb{Q}}[H]$, has a direct financial interpretation
grounded in martingale duality: it equals the minimal initial capital
required to construct an adapted, non-negative wealth process that
dominates~$H$ almost surely at maturity~$T$.  This is the
{minimal dominating price}.  In a complete market, the interval
degenerates to a single point and the dominating strategy reduces to the
familiar replicating portfolio with no residual cost.

American options introduce an additional layer of complexity through
optimal stopping.  Unlike European claims, which are exercised only at
maturity, an American option may be exercised by the holder at any stopping
time $\tau \le T$.  The seller's hedge must remain adequate regardless of
when the holder exercises.  This couples the hedging problem with an
optimal stopping problem: the value to the seller must account
simultaneously for the worst-case choice of pricing measure {and}
the worst-case exercise policy.  The minimal dominating price becomes a
nested essential supremum over both sources of adversarial choice:
\begin{equation}
\label{eq:american_price_intro}
\pi(\xi)
\;=\;
\sup_{\mathbb{Q} \in \mathcal{M}(X)}
\;\sup_{\tau \in \mathcal{T}_{[0,T]}}
\;\mathbb{E}^{\mathbb{Q}}[\xi_\tau],
\end{equation}
where $\xi = (\xi_t)_{t \in [0,T]}$ is the discounted reward process of
the American option and $\mathcal{T}_{[0,T]}$ is the set of stopping times
valued in $[0,T]$.  In complete markets, the outer supremum is trivial and
the problem reduces to the classical Snell envelope construction of
\citet{Bensoussan1984} and \citet{Karatzas1988}.  In incomplete markets,
both suprema are non-trivial and their interaction creates a substantially
harder problem with a minimax structure.

The central object of this thesis is the {minimal dominating wealth
process} under non-unique martingale measures: the adapted process that
tracks the value~\eqref{eq:american_price_intro} dynamically through time,
together with the trading strategy that generates it.  The thesis does not
introduce new pricing theory or new representation results; rather, it
investigates this object computationally using established deep learning
methods, and interprets the results through the lens of existing
martingale duality results from mathematical finance.  The non-uniqueness
of martingale measures creates a pricing interval whose upper bound defines
the minimal dominating price, and American options compound the difficulty
through the interaction of model ambiguity with optimal stopping.  The
precise mathematical formulation of the problem is the subject of the next
section.


% =========================================================================
% 1.2  MATHEMATICAL FORMULATION
% =========================================================================
\subsection{Mathematical Formulation of the Problem}
\label{subsec:math_formulation}

Let $(\Omega, \mathcal{F},
\mathbb{F} = (\mathcal{F}_t)_{t \ge 0}, \mathbb{P})$ be a filtered
probability space satisfying the usual conditions of right-continuity and
completeness.  The space supports $d = 2$ traded risky assets with
discounted price process $X := (X^1, X^2)$, assumed to be a locally
bounded, càdlàg (right-continuous with left limits) semimartingale.  Let
$\mathcal{M}(X)$ denote the set of all probability measures
$\mathbb{Q} \sim \mathbb{P}$ under which $X$ is a local martingale.  By
the first fundamental theorem of asset pricing
\citep{DelbaenSchachermayer1994}, the absence of arbitrage (NFLVR) is
equivalent to $\mathcal{M}(X) \neq \emptyset$.  The market is assumed to
be incomplete: $|\mathcal{M}(X)| > 1$, reflecting the presence of
$m > d = 2$ independent sources of randomness that cannot all be hedged by
trading in the two assets.

A {trading strategy} is a predictable, $X$-integrable process
$\varphi := (\varphi^1, \varphi^2)$ representing the number of units of
each asset held at each time.  Given an initial capital $v \ge 0$ and a
trading strategy~$\varphi$, the {wealth process} is
\begin{equation}
\label{eq:wealth_process}
V_t
\;:=\; v \;+\; \int_0^t \varphi_s \, dX_s \;-\; C_t,
\qquad t \ge 0,
\end{equation}
where $\int_0^t \varphi_s \, dX_s$ represents the cumulative gains from
trade and $C = (C_t)_{t \ge 0}$ is an adapted, non-decreasing process
with $C_0 = 0$.  When $C \equiv 0$, the strategy is {self-financing}:
changes in wealth are driven solely by changes in asset prices.  When
$C$ is strictly increasing, the strategy requires the hedger to absorb a
cumulative cost that arises because the market does not offer enough
instruments to track the obligation perfectly.  Following \citet{Kramkov1996},
the strategy is called {admissible} if $V_t \ge 0$ for all $t \ge 0$.

The American put option with strike~$K$ and maturity~$T$ on asset~$i$
generates a discounted reward process
\begin{equation}
\label{eq:american_reward}
\xi_t \;:=\; e^{-rt}\max(K - S^i_t,\, 0),
\qquad t \in [0, T],
\end{equation}
that the holder may exercise at any stopping time $\tau \le T$.  An
admissible strategy with wealth $V_t \ge \xi_t$ for all
$t \in [0,T]$ is called a {dominating strategy} for~$\xi$.  The
{minimal dominating value process} is
\begin{equation}
\label{eq:value_process}
\tilde{V}_t
\;:=\;
\esssup_{\mathbb{Q} \in \mathcal{M}(X)}
\;\esssup_{\tau \in \mathcal{T}_{[t,T]}}
\;\mathbb{E}^{\mathbb{Q}}[\xi_\tau \mid \mathcal{F}_t],
\end{equation}
where $\mathcal{T}_{[t,T]}$ is the set of $\mathbb{F}$-stopping times
valued in $[t,T]$.  At time zero, $\tilde{V}_0$ reduces to the minimal
dominating price~\eqref{eq:american_price_intro}.

Under appropriate integrability conditions, the {optional
decomposition theorem} of \citet{Kramkov1996} guarantees that the process
$\tilde{V}$ admits the representation
\begin{equation}
\label{eq:optional_decomp}
\tilde{V}_t
\;=\; \tilde{V}_0
\;+\; \int_0^t \tilde{\varphi}_s \, dX_s
\;-\; \tilde{C}_t,
\qquad t \ge 0,
\end{equation}
where $\tilde{\varphi}$ is a predictable, $X$-integrable process and
$\tilde{C}$ is an adapted, non-decreasing, càdlàg process with
$\tilde{C}_0 = 0$.  The decomposition is called {optional} because
$\tilde{C}$ is $\mathcal{F}_t$-measurable at each time~$t$ but not
necessarily $\mathcal{F}_{t^-}$-measurable, in contrast to the classical
Doob--Meyer decomposition in which the increasing process is
{predictable} \citep{KaratzasShreve1991}.  In a complete market,
$\tilde{C} \equiv 0$ and the decomposition reduces to a self-financing
replication.  In an incomplete market, $\tilde{C}$ is strictly increasing
and measures the irreducible cost of the gap between the instruments
available and the risks to be managed.  The full statement and financial
interpretation of the optional decomposition theorem are presented in
Chapter~2.

The theorem is {existential}.  It guarantees the existence of the
predictable integrand~$\tilde{\varphi}$ and the adapted cost
process~$\tilde{C}$, but it does not provide a constructive algorithm for
computing either object.  In any concrete setting with stochastic
volatility, multiple assets, or early-exercise features, one must resort
to numerical methods.  This gap between the existence guarantee and the
need for a computable solution defines the problem addressed in this
thesis: to approximate the integrand~$\tilde{\varphi}$ using established
neural network methods, and to interpret the residual hedging error as an
empirical estimate of the cost process~$\tilde{C}$.


% =========================================================================
% 1.3  COMPUTATIONAL CHALLENGE
% =========================================================================
\subsection{Computational Challenge}
\label{subsec:computational_challenge}

The optional decomposition~\eqref{eq:optional_decomp} identifies the
structural target, a predictable integrand and an adapted cost process,
but its computation in concrete settings requires solving
high-dimensional, path-dependent optimisation problems for which no
closed-form solutions are available.  This section reviews the principal
classes of numerical methods, explains their limitations, and motivates
the adoption of established deep learning approaches.

\subsubsection{Classical PDE Methods}

When the price process is Markovian and low-dimensional, the value
process~\eqref{eq:value_process} satisfies a partial differential equation
(PDE) or, for American options, a free-boundary problem.  For the
Black--Scholes model with a single asset, the American put price satisfies
a well-studied complementarity problem that can be solved by finite
differences or front-fixing methods \citep{Wilmott1995}.  However, these
methods scale poorly with the number of state variables.  In the two-asset
Heston model considered in this thesis, the state space is at least
four-dimensional (two asset prices and two variance processes), and the
PDE becomes a system of coupled equations on a high-dimensional grid.
Storage and computation costs grow exponentially with dimension, the
well-known ``curse of dimensionality,'' making grid-based methods
impractical for all but the simplest multi-asset settings
\citep{Glasserman2003}.

Moreover, PDE methods yield the value function and its derivatives (the
Greeks), from which the hedging strategy can be extracted.  But when the
market is incomplete, the value function depends on the choice of pricing
measure, and no single PDE characterises the minimal dominating value
process without additional information about the structure
of~$\mathcal{M}(X)$.  This model-dependence is a fundamental limitation
that persists even when the dimensionality is manageable.

\subsubsection{Regression-Based Monte Carlo Methods}

Monte Carlo simulation handles high-dimensional problems naturally, since
its convergence rate is independent of the number of state variables.  For
American options, the key difficulty is that the optimal exercise boundary
is not known in advance and must be estimated from the simulated paths
themselves.  The {Longstaff--Schwartz algorithm}
\citep{LongstaffSchwartz2001} addresses this by regressing the
continuation value onto a set of basis functions of the current state at
each time step, proceeding backwards from maturity.  The
{Tsitsiklis--Van~Roy} algorithm \citep{TsitsiklisVanRoy2001}
provides a closely related approach with convergence guarantees.

These regression-based methods are widely used in practice and serve as
important baselines.  Their principal limitation is the choice of basis
functions: in high-dimensional or path-dependent settings, it is difficult
to select a set of regressors that is both computationally tractable and
expressive enough to capture the continuation value accurately.
Furthermore, the regression step provides an approximation of the
{optimal stopping rule} but does not directly yield an optimal
{hedging strategy}: a separate computation is needed to determine
the portfolio positions along each path.  In incomplete markets, the
additional challenge of accounting for model ambiguity across
$\mathcal{M}(X)$ is not naturally addressed by these methods.

\subsubsection{Deep Hedging}

The {deep hedging} approach of
\citet{BuehlerGononTeichmannWood2019} takes a fundamentally different
route.  Rather than solving for the value function and then extracting a
strategy, it parameterises the trading strategy directly as a neural
network and trains it end-to-end by minimising a risk functional applied to
the distribution of hedging errors.  At each rebalancing date, the network
maps the current market state (asset prices, time to maturity, and any
auxiliary features) to a portfolio position.  A large batch of price
paths is simulated, the strategy is applied along each path, and the
terminal hedging error is evaluated.  The network weights are updated by
stochastic gradient descent through the entire trading sequence.

Deep hedging is model-agnostic: it does not require analytic pricing
formulas, and transaction costs, position limits, and non-standard payoffs
can be incorporated within the same optimisation.
\citet{BuehlerGononTeichmannWood2019} proposed two primary architectures:
a feedforward network applied at each time step, and a recurrent network
that maintains a hidden state encoding the path history.
\citet{Horvath2021} introduced path signatures as input features,
replacing hand-crafted features with a systematic mathematical
representation of path information.  The literature review in Chapter~2
surveys the subsequent extensions in detail, including transformer
architectures \citep{Cao2023} and reinforcement learning formulations
\citep{CaoChenVorobeychikMadan2023}.

\subsubsection{Deep BSDE Solvers}

An alternative computational approach, and the one that serves as the
{base model} in this thesis, is the {deep BSDE solver} of
\citet{HanJentzenE2018}.  Backward stochastic differential equations
(BSDEs), introduced in the non-linear setting by
\citet{PardouxPeng1990} and applied to finance by
\citet{ElKarouiPengQuenez1997}, provide a natural mathematical
formulation of hedging problems.  In the BSDE setting, the portfolio
value~$Y_t$ and the hedge ratio~$Z_t$ satisfy a coupled
forward--backward system: the forward SDE describes the evolution of the
underlying prices, while the backward equation propagates the terminal
condition (the option payoff) back through time.  The connection to the
optional decomposition is direct: the process~$Y$ corresponds to the
wealth process~$V$, the control~$Z$ corresponds to the
integrand~$\tilde{\varphi}$, and the driver of the backward equation
encodes the cost structure associated with~$\tilde{C}$.

The deep BSDE solver approximates the solution by parameterising the
initial value~$Y_0$ and the control process~$Z_{t_k}$ at each discrete
time step~$t_k$ with feedforward sub-networks, and training end-to-end by
minimising the mean squared terminal error.  \citet{HanJentzenE2018}
demonstrated that this approach scales to problems in hundreds of
dimensions where classical PDE-based methods are infeasible.
\citet{HurePhamWarin2020} extended the methodology to control problems
and \citet{Gobet2023} applied it to hedging under incomplete markets.
The deep BSDE solver is adopted as the base model in this thesis because
of its direct structural correspondence with the backward equation that
characterises the minimal dominating value process.

\subsubsection{Path Generation via Variational Autoencoders}

Training any of the above deep learning methods requires large numbers of
simulated price paths.  In addition to standard Monte Carlo simulation
from parametric models, this thesis employs a {variational
autoencoder} (VAE) \citep{KingmaWelling2014} trained on historical
multi-asset equity data to generate synthetic sample paths.  The VAE
learns a low-dimensional latent representation of the joint distribution
of asset returns and produces synthetic trajectories that capture empirical
features, such as volatility clustering, fat tails, and cross-asset
dependence structures, that parametric models may not fully reproduce.
This data-driven path generation follows the approach advocated by
\citet{Wiese2020} and \citet{Ni2021}, and provides a robustness check
against model misspecification: by evaluating hedging performance on both
parametric and VAE-generated paths, one can assess the extent to which
conclusions depend on the assumed data-generating process.  The historical
dataset used and the VAE architecture are described in Chapter~3.

\subsubsection{Scope and Loss Function}

The thesis does not derive or extend any of the methods reviewed above.
The deep BSDE solver, the deep hedging architectures, and the VAE path
generator are all established methods with published convergence results
and documented implementations.  The contribution is computational and
empirical: implementing these methods under identical experimental
conditions, applying them to the specific problem of hedging American put
options in controlled incomplete-market settings, and comparing their
performance.

All three hedging architectures (deep BSDE, GRU, and FNN with path
signatures) are trained using a common asymmetric loss function that
penalises under-hedging more heavily than over-hedging:
\begin{equation}
\label{eq:loss}
\mathcal{L}(\theta)
\;:=\; \lambda_s \, \mathbb{E}\bigl[\max(\xi_\tau - V_T^\theta,\, 0)\bigr]
\;+\; \lambda_o \, \mathbb{E}\bigl[\max(V_T^\theta - \xi_\tau,\, 0)\bigr]
\;+\; \lambda_c \, \operatorname{CVaR}_\alpha\!\bigl(\max(\xi_\tau
      - V_T^\theta,\, 0)\bigr),
\end{equation}
where $V_T^\theta$ is the terminal portfolio value under strategy
$\varphi_\theta$, and $\operatorname{CVaR}_\alpha$ denotes the conditional
value-at-risk at level~$\alpha$ \citep{RockafellarUryasev2000}.  The
asymmetry $\lambda_s \gg \lambda_o$ reflects the practical reality that
failing to cover an obligation is financially dangerous, whereas holding a
modest surplus is merely capital-inefficient.  The CVaR term targets the
tail of the shortfall distribution, penalising not just the average
shortfall but the severity of the worst outcomes.  This loss function does
not enforce the strict domination $V_t \ge \xi_t$ for all~$t$ that would
characterise a full super-replicating strategy; instead, it encodes a
{soft} preference for strategies whose shortfall is small and whose
tail risk is controlled.  From the perspective of the optional
decomposition, the mean residual hedging error of the trained model
provides an empirical estimate of the expected value of the cost
process~$\tilde{C}$, while the shortfall distribution characterises the
pathwise behaviour of~$\tilde{C}$ across scenarios.  The precise
calibration of the penalty parameters $\lambda_s$, $\lambda_o$, and
$\lambda_c$ is described in Chapter~4.

The problem of computing the minimal dominating wealth process and its
associated trading strategy is therefore computationally challenging:
classical methods face scalability and model-dependence limitations, and
established deep learning methods offer viable alternatives.  The thesis
performs a comparative empirical study of these alternatives, as specified
by the research questions in the next section.


% =========================================================================
% 1.4  RESEARCH QUESTIONS
% =========================================================================
\subsection{Research Questions}
\label{subsec:research_questions}

The thesis addresses the following empirical research questions.

\begin{enumerate}
\item \textbf{Approximation of the minimal dominating wealth process.}
  Can neural-network-based trading strategies approximate the minimal
  dominating wealth process implied by martingale duality in multi-asset
  incomplete markets?  Specifically, how closely do the portfolio value
  processes generated by the deep BSDE solver, the GRU network, and the
  FNN with path signature features track the value
  process~\eqref{eq:value_process}, and how does each architecture's
  residual error compare with the others under identical market
  conditions?

\item \textbf{Structural versus dynamic incompleteness.}
  How does the source of market incompleteness affect the magnitude and
  distribution of the residual hedging error?  The multi-factor GBM model
  introduces incompleteness through a dimensional mismatch between traded
  assets and Brownian drivers (a structural feature fixed at time zero),
  while the Heston model introduces it through the random evolution of
  untraded variance processes (a dynamic feature that changes over time).
  Comparing the two settings isolates the contribution of each mechanism
  to the total hedging cost.

\item \textbf{Recurrent memory and adapted cost components.}
  The cost process~$\tilde{C}$ in the optional
  decomposition~\eqref{eq:optional_decomp} is adapted but generally not
  predictable \citep{Kramkov1996}: its value at time~$t$ may depend on
  information revealed at~$t$ itself, not merely on information available
  before~$t$.  The GRU architecture maintains an evolving hidden state
  that is updated at each observation, while the FNN with path signatures
  constructs a static, fixed-dimensional summary of the path history.
  Does recurrent memory improve the approximation of adapted cost
  components relative to feedforward or signature-based architectures?

\item \textbf{Deep BSDE versus direct deep hedging.}
  The deep BSDE solver approximates the hedging problem through a
  backward equation, while the GRU and FNN-Sig architectures parameterise
  the strategy directly and optimise a forward risk functional.  How do
  these two approaches compare in terms of mean hedging error, tail risk
  (measured by CVaR), and computational cost?

\item \textbf{Robustness to model misspecification.}
  How do hedging models trained on VAE-generated paths, which reflect
  empirical distributional features learned from historical data, compare
  with those trained on parametric (GBM, Heston) simulations?  Does
  data-driven path generation improve out-of-sample hedging performance
  and reduce sensitivity to model assumptions?
\end{enumerate}

Each question is empirical in nature.  The thesis does not seek to prove
new theoretical results, but to provide computational evidence on the
performance of established methods in a setting, multi-asset American
options under controlled incomplete-market models, that has not been
studied in this specific combination.  The questions are designed so that
the answers depend on measurable quantities (mean hedging error, shortfall
probability, CVaR, and the distribution of the residual cost) that can be
estimated from simulated experiments with controlled statistical
uncertainty.


% =========================================================================
% 1.5  CONTRIBUTIONS
% =========================================================================
\subsection{Contributions}
\label{subsec:contributions}

The contributions of this thesis are computational and empirical.  All
mathematical results, pricing models, representation theorems, and
numerical solvers are drawn from the established literature.

\begin{enumerate}
\item \textbf{Comparative evaluation under identical conditions.}
  The thesis provides a systematic empirical comparison of three
  neural-network-based hedging architectures: the deep BSDE solver
  \citep{HanJentzenE2018}, a GRU recurrent network
  \citep{BuehlerGononTeichmannWood2019, Cho2014}, and a feedforward
  network with path signature features \citep{Horvath2021}.  All three
  are applied to the same hedging problem under the same market models,
  loss function, and evaluation protocol.  A two-stage experimental
  design, consisting of Bayesian hyperparameter optimisation
  \citep{Snoek2012} followed by retraining across multiple random seeds
  \citep{Henderson2018, Bouthillier2021}, ensures that reported
  performance reflects genuine model capability rather than fortunate
  initialisation or hyperparameter selection.

\item \textbf{Two-model study isolating sources of incompleteness.}
  By comparing a multi-driver GBM model (structural incompleteness)
  with a two-asset Heston model (dynamic incompleteness), the thesis
  disentangles the contributions of dimensional mismatch and
  stochastic volatility to the residual hedging error.  This provides
  empirical evidence on how the {source} of incompleteness, not
  merely its presence, affects hedging difficulty.

\item \textbf{Tail-risk-based evaluation.}
  Performance is evaluated using the asymmetric loss
  function~\eqref{eq:loss} with a CVaR component, providing a
  risk-sensitive assessment that goes beyond mean hedging error to
  characterise the tail of the shortfall distribution.  This evaluation
  is more informative than mean-squared-error comparisons for
  applications in which the cost of under-hedging is asymmetric.

\item \textbf{Empirical proxy for the cost of incompleteness.}
  The residual hedging error of the trained models is interpreted as
  an empirical estimate of the adapted cost process~$\tilde{C}$ in
  the optional decomposition~\eqref{eq:optional_decomp}.  The thesis
  provides empirical evidence consistent with the theoretical
  prediction that this cost is strictly positive in incomplete markets
  and approximately zero in the complete-market limit.

\item \textbf{VAE-based path generation and robustness analysis.}
  The use of a variational autoencoder trained on historical data
  provides a data-driven complement to parametric simulation and
  enables an assessment of model misspecification effects on hedging
  performance.
\end{enumerate}

The thesis does not extend or modify any of the theoretical results or
numerical methods cited above.  It uses established methods as
computational tools and evaluates them on a problem that, to the best of
the author's knowledge, has not been studied in this specific combination
of models, instruments, and architectures.  Results are described as
providing empirical evidence consistent with theoretical predictions, not
as extending or proving those predictions.


% =========================================================================
% 1.6  STRUCTURE OF THE THESIS
% =========================================================================
\subsection{Structure of the Thesis}
\label{subsec:thesis_outline}


